\section{Signed-inverse families at prime-power steps}
\label{sec:signed-inverse}

For every prime power $Q=p^e$, this section defines a finite family $\mathcal A_Q\subset\mathcal C$ and its simple-quotient map $\phi_Q:\mathcal A_Q\to S_Q^\times$.  The construction is allowed to be empty at small $Q$; the density, fibre, incompatibility, support and reciprocal-mass estimates used later are eventual.  A single onset $Q_{\mathrm{si}}$ for all of those later capabilities is fixed at the end of the section.

Let $\Lambda$ be the integer supplied by Proposition~\ref{prop:prime-supply-main}, and put
\[
B_{\mathrm{car}}=\Lambda+1.
\]
Let $Q=p^e$ and define
\begin{equation}\label{eq:carrier-base}
X_Q=\left\lfloor\frac{Q}{B_{\mathrm{car}}}\right\rfloor.
\end{equation}
Define the carrier set
\begin{equation}\label{eq:carrier-band}
\mathscr B_Q=
\{b\text{ prime}:X_Q<b\le\Lambda X_Q,\ (b,Q)=1\}.
\end{equation}
Writing
\[
Q=B_{\mathrm{car}}X_Q+r,
\qquad 0\le r<B_{\mathrm{car}},
\]
shows that every $b\in\mathscr B_Q$ satisfies
\begin{equation}\label{eq:carrier-exact}
B_{\mathrm{car}}b>Q,
\qquad
b\le\Lambda X_Q<B_{\mathrm{car}}X_Q\le Q.
\end{equation}
For $X_Q\ge1$ one also has
\[
B_{\mathrm{car}}X_Q\le Q<2B_{\mathrm{car}}X_Q,
\]
so $X_Q\asymp Q$.  Once $X_Q$ lies beyond the onset in Proposition~\ref{prop:prime-supply-main}, that proposition gives $\asymp Q/\log Q$ possible carriers; imposing $(b,Q)=1$ removes at most the prime $p$.

For each $b\in\mathscr B_Q$, let $c_+\in\{1,\ldots,Q-1\}$ be the inverse of $b$ modulo $Q$, put $c_-=Q-c_+$, and write
\begin{equation}\label{eq:signed-inverse}
bc_+=Qk_++1,
\qquad
bc_-=Qk_--1.
\end{equation}
Then
\[
1\le k_\pm<b<Q,
\qquad
k_++k_-=b.
\]
Since $p\nmid b$, the two coefficients cannot both be divisible by $p$.  Define a sign deterministically: if $e=1$, use the sign with larger coefficient, breaking a tie in favour of $+$; if $e\ge2$, use $+$ unless $p\mid k_+$, and use $-$ otherwise.  Let $k=k(Q,b)$ be the coefficient determined by this rule, and define
\begin{equation}\label{eq:signed-block}
I(Q,b)=
\begin{cases}
[Qk,Qk+1],&+,\\
[Qk-1,Qk],&-.
\end{cases}
\end{equation}
If $Q=p$ is prime, this rule gives
\begin{equation}\label{eq:atomic-k-lower}
k\ge b/2>X_Q/2\gg p.
\end{equation}

\subsection{Transversality and deterministic deduplication}

The distinguished denominator $Qk$ has largest prime-power stage exactly $Q$: the coefficient $k$ is a $p$-unit, and every prime-power factor contributed by $k$ has rank below $Q$ because $k<Q$.

\begin{lemma}[Carrier transversality and coefficient multiplicity]\label{lem:carrier-transversality}
For all sufficiently large prime powers $Q=p^e$ the following hold.
\begin{enumerate}
\item All but $O_{B_{\mathrm{car}}}(1)$ carrier primes have companion denominator with no prime-power factor of rank at least $Q$.
\item A fixed coefficient $k$ is produced by at most four carrier primes.
\end{enumerate}
\end{lemma}

\begin{proof}
For the companion denominator $bc_\pm=Qk\pm1$, both factors $b$ and $c_\pm$ are below $Q$.  Suppose that a prime power $\ell^r\ge Q$ divides $bc_\pm$.  If $\ell\ne b$, then the primality of $b$ forces $\ell^r\mid c_\pm$, contradicting $c_\pm<Q$.  Hence $\ell=b$.  Since also $b<Q$, one has $r\ge2$, and therefore $b\mid c_\pm$.  Write $c_\pm=\ell_0 b$.  Since $c_\pm<Q<B_{\mathrm{car}}b$, one has $1\le\ell_0<B_{\mathrm{car}}$, and \eqref{eq:signed-inverse} gives
\begin{equation}\label{eq:bad-carrier}
\ell_0 b^2\equiv\pm1\pmod Q.
\end{equation}
For any $(\ell_0,\pm)$, if $p\mid\ell_0$, the congruence has no solution because its right-hand side is a unit modulo $p$.  Otherwise it reduces to a unit square congruence modulo $p^e$, which has at most two roots for odd $p$ and at most four roots for $p=2$.  For the latter assertion, when $e\ge3$ the quotient of any two unit roots is a solution of $u^2\equiv1\pmod{2^e}$, of which there are four; the cases $e\le2$ are immediate.  There are only $2(B_{\mathrm{car}}-1)$ choices of $(\ell_0,\pm)$, proving the first assertion.

For any coefficient $k$, every carrier producing it divides one of $Qk-1$ and $Qk+1$.  Since $1\le k<Q$, both integers are smaller than $Q^2$.  Every carrier satisfies $b>Q/B_{\mathrm{car}}$ by \eqref{eq:carrier-base} and \eqref{eq:carrier-band}.  If $Q>B_{\mathrm{car}}^3$, a number below $Q^2$ cannot contain three distinct carrier primes, because their product would exceed $Q^3/B_{\mathrm{car}}^3>Q^2$.  Thus each of $Qk-1$ and $Qk+1$ contains at most two possible carriers, giving multiplicity at most four.
\end{proof}

For every $Q$, discard exactly those carriers whose companion denominator has a prime-power factor of rank at least $Q$.  Lemma~\ref{lem:carrier-transversality} shows that this deletes only $O_{B_{\mathrm{car}}}(1)$ carriers once $Q$ is sufficiently large.  Different remaining carriers can still produce the same coefficient $k$; keep, for each coefficient, the occurrence with least carrier.  The eventual multiplicity bound in the lemma shows that, for sufficiently large $Q$, the resulting coefficient set has cardinality
\[
\asymp Q/\log Q.
\]

If $e\ge2$, order the remaining coefficients $k_1<\cdots<k_N$ and retain the upper half.  Since $k_i\ge i$, every retained coefficient then satisfies
\begin{equation}\label{eq:coef-lower}
k\gg Q/\log Q.
\end{equation}
If $Q=p$ is prime, retain all coefficients.  Let $\mathcal A_Q$ be the resulting family of intervals \eqref{eq:signed-block}.

\begin{lemma}[Top-stage recovery]\label{lem:top-stage-recovery}
For every $I\in\mathcal A_Q$, exactly one of its two denominators has largest prime-power stage $Q$, namely the distinguished denominator $Qk$.  Consequently $I$ determines its stage $Q$ and coefficient $k$; in particular the families $\mathcal A_Q$ for distinct $Q$ are disjoint.
\end{lemma}

\begin{proof}
The distinguished denominator has top stage $Q$ because $k$ is a $p$-unit and $k<Q$.  The transversality deletion removes every carrier for which the companion denominator has a prime-power factor of rank at least $Q$.  Hence the companion has strictly smaller top stage.
\end{proof}

\subsection{The simple-quotient value}

Modulo $F_{<Q}$, the companion reciprocal has only smaller prime-power stages, so the class contributed by $I(Q,b)$ is the class of $1/(Qk)$.  Let $u\in\{1,\ldots,p-1\}$ be the unique inverse of $k$ modulo $p$, and write
\[
uk=1+pt.
\]
Then
\begin{equation}\label{eq:simple-coordinate-calculation}
\frac{u}{Q}-\frac1{Qk}
=\frac{t}{p^{e-1}k}.
\end{equation}
Every prime-power factor of the denominator on the right has rank below $Q$: the $p$-part is at most $p^{e-1}$, and every factor coming from $k$ is below $Q$.  Hence the difference belongs to $F_{<Q}$.  In particular,
\begin{equation}\label{eq:signed-label-filtration}
\overline W(I(Q,b))\in F_Q,
\end{equation}
and, relative to the distinguished generator represented by $1/Q$, its image in $S_Q=F_Q/F_{<Q}$ is
\begin{equation}\label{eq:simple-value}
\phi_Q(I(Q,b))=k^{-1}\pmod p.
\end{equation}
Thus the simple-quotient map
\[
\phi_Q:\mathcal A_Q\longrightarrow S_Q^\times
\]
is well defined.  By Lemma~\ref{lem:top-stage-recovery}, the distinguished-point map $I(Q,b)\mapsto Qk$ is globally injective.

For an integer interval $J=[u,v]\subset\N_{>0}$ define its \emph{exclusion envelope}
\[
\operatorname{Env}(J)
=\{n\in\N_{>0}:\operatorname{dist}(n,J)\le1\}.
\]
Thus another interval component is incompatible with $J$ exactly when its support meets $\operatorname{Env}(J)$.  A member of some $\mathcal A_Q$ is said to meet an exclusion envelope if its support intersects that envelope.

\begin{proposition}[Uniform estimates for the signed-inverse families]\label{prop:signed-inverse-capabilities}
The deterministic families $\mathcal A_Q$ satisfy the following properties.
\begin{enumerate}
\item Their normalized density has a positive lower limit:
\begin{equation}\label{eq:family-density-liminf}
\kappa_\infty
:=
\liminf_{\substack{Q\to\infty\\Q\text{ prime power}}}
\frac{|\mathcal A_Q|\log Q}{Q}
>0.
\end{equation}
\item For every sufficiently large $Q=p^e$ and every $s\in S_Q^\times$,
\begin{equation}\label{eq:fibre-bound}
|\phi_Q^{-1}(s)|\le Q/p.
\end{equation}
\item The incompatibility graph on $\bigcup_Q\mathcal A_Q$ has finite maximum degree.  More generally, for every fixed $L\ge1$ there is $C_L$ such that an exclusion envelope of cardinality at most $L$ meets at most $C_L$ members of the union.
\item The normalized lower support edge has positive lower limit:
\begin{equation}\label{eq:support-liminf}
c_\infty
:=
\liminf_{\substack{Q\to\infty\\Q\text{ prime power}}}
\left(
\min_{\substack{I\in\mathcal A_Q\\z\in I}}
\frac{z\log Q}{Q^2}
\right)
>0,
\end{equation}
where the minimum is taken only for the eventually nonempty families.  Every support point also satisfies $z\le2Q^2$ for all sufficiently large $Q$, and if $Q$ is prime then $z\gg Q^2$.
\item Define
\begin{equation}\label{eq:canonical-role-mass}
\omega_Q
=
\max\left(
\{0\}\cup
\left\{
W\!\left(\bigcup_{I\in O}I\right):
O\subseteq\mathcal A_Q,\ |O|\le p,\ \bigcup_{I\in O}I\in\mathcal C
\right\}
\right).
\end{equation}
Then
\begin{equation}\label{eq:role-summable}
\sum_Q\omega_Q<\infty.
\end{equation}
\end{enumerate}
\end{proposition}

\begin{proof}
The positive lower density follows from the carrier count, the bounded exceptional set, bounded coefficient multiplicity and, for $e\ge2$, the upper-half thinning.  For \eqref{eq:fibre-bound}, equality of simple-quotient values implies equality of $k$ modulo $p$.  A fixed nonzero residue class modulo $p$ occurs exactly $Q/p$ times among $1\le k<Q=p^e$, and the deterministic deduplication retains at most one interval for each coefficient.

For incompatibility, each interval occupies two adjacent integers around its globally injective distinguished point $Qk$.  If an exclusion envelope has cardinality at most $L$, every member meeting it has distinguished point in a fixed one-step enlargement of the envelope.  This enlargement contains $O(L)$ integer sites and global injectivity leaves at most one retained interval per site.  In particular the global incompatibility graph has bounded integer degree.

The positive lower support limit follows from $z=Qk+O(1)$ together with \eqref{eq:coef-lower}; the prime lower bound follows from \eqref{eq:atomic-k-lower}.  The upper bound $z\le2Q^2$ follows from $k<Q$.

For the quantity \eqref{eq:canonical-role-mass}, a prime family has $O(p/\log p)$ intervals, each of reciprocal mass $O(p^{-2})$, so
\[
\omega_p=O\!\left(\frac1{p\log p}\right).
\]
For $Q=p^e$, $e\ge2$, each interval has mass $O(\log Q/Q^2)$, so
\[
\omega_{p^e}
=O\!\left(\frac{p\log Q}{Q^2}\right)
=O\!\left(\frac{e\log p}{p^{2e-1}}\right).
\]
For the prime terms, let $X_j=\Lambda^jX_0$, with $X_0$ beyond the onset of Proposition~\ref{prop:prime-supply-main}.  The contribution of primes in $X_j<p\le\Lambda X_j$ is at most
\[
O\!\left(\frac{X_j}{\log X_j}\right)
O\!\left(\frac1{X_j\log X_j}\right)
=O\!\left(\frac1{(\log X_j)^2}\right)
=O(j^{-2}),
\]
so the prime terms in the tail are summable.  The composite tail is bounded by a constant multiple of $\sum_p p^{-2}$ after summing over $e\ge2$.  The finitely many remaining prime-power stages contribute a finite amount, hence \eqref{eq:role-summable} holds.
\end{proof}

Set the positive constants
\begin{equation}\label{eq:canonical-density-support-constants}
\kappa=\frac{\kappa_\infty}{2},
\qquad
c=\frac12\min\{1,c_\infty\}.
\end{equation}
Here the lower limits are allowed to take values in the extended real line; the truncation ensures that $c$ is finite even if $c_\infty=+\infty$.  The carrier upper bound makes $\kappa_\infty$ finite.  By the definitions of the lower limits, for every sufficiently large prime power $Q$ one has
\begin{equation}\label{eq:family-size}
|\mathcal A_Q|\ge\kappa\frac Q{\log Q}
\end{equation}
and every support point satisfies
\begin{equation}\label{eq:support-window}
c\frac{Q^2}{\log Q}\le z\le2Q^2.
\end{equation}

Choose once and for all an integer $Q_{\mathrm{si}}$ so large that, for every prime power $Q\ge Q_{\mathrm{si}}$, the carrier-transversality and coefficient-multiplicity bounds, the fibre bound \eqref{eq:fibre-bound}, the family-size estimate \eqref{eq:family-size}, and the support bounds \eqref{eq:support-window} (including the prime lower bound used above) all hold.  The objects $\mathcal A_Q$, $\phi_Q$ and $\omega_Q$ remain defined for every prime power $Q$; the finitely many stages below $Q_{\mathrm{si}}$ are simply outside the later transport regime.  Every transport onset introduced below is taken to be at least $Q_{\mathrm{si}}$.

Let $G_{\infty}$ be the incompatibility graph on $\bigcup_Q\mathcal A_Q$.  Proposition~\ref{prop:signed-inverse-capabilities} shows that the bounded subset
\[
\{\deg_{G_{\infty}}(v):v\in V(G_{\infty})\}\subseteq\N
\]
has a maximum; define
\begin{equation}\label{eq:canonical-global-degree}
\Delta_*=\max_{v\in V(G_{\infty})}\deg_{G_{\infty}}(v).
\end{equation}
Abbreviate
\begin{equation}\label{eq:packing-constants}
L_P=L_P(\Delta_*),
\qquad
\rho_P=\rho_P(\Delta_*),
\end{equation}
with the functions of \eqref{eq:haxell-canonical-parameters}.  These constants will be used in the finite compatibility step.

In particular, for every $\varepsilon>0$ there is $Q_1$ such that every finite set of prime-power steps of rank at least $Q_1$ has total reciprocal-mass bound below $\varepsilon$.
